\documentclass{article}
\usepackage[margin=1.15in]{geometry}
\usepackage{xcolor}
\usepackage{parskip}
\usepackage{fancyhdr}
\usepackage{array}
\usepackage{booktabs}
\usepackage{enumitem}
\usepackage{amsmath}
\usepackage{tabularx}
\usepackage{graphicx}

\pagestyle{fancy}
\fancyhf{}
\cfoot{\thepage}
\rhead{Lecture 3}
\lhead{MA 214: Applied Statistics}
\renewcommand{\headrulewidth}{.4mm}

\newcommand{\blankline}[1]{\underline{\hspace{#1}}}

\begin{document}

\noindent\textbf{Name:}\blankline{2.25in}\hfill\textbf{BUID:}\blankline{1.55in}\newline

\begin{center}
 \Large In-Class Worksheet: Simple Linear Regression
\end{center}

\subsection*{Scenario}
Simple linear regression models a numerical response using one numerical explanatory variable. Today we practice identifying variable roles, interpreting fitted lines, computing residuals, and connecting least squares to the regression model.

\medskip
\noindent\textbf{Goals:} \textbf{(1)} identify response and explanatory variables, \textbf{(2)} state basic model assumptions, \textbf{(3)} interpret slope and intercept, \textbf{(4)} compute residuals and RSS.

\subsection*{Part 1: Response and Explanatory Variables}
For each research question, identify the explanatory variable and the response variable. Remember: these roles depend on the research question and do not automatically imply causation.

\begin{center}
\renewcommand{\arraystretch}{1.45}
\begin{tabularx}{\textwidth}{p{7.3cm}p{4cm}X}
\toprule
\textbf{Research question} & \textbf{Explanatory variable} & \textbf{Response variable} \\
\midrule
Can commute distance help predict commute time? & & \\[1.8em]
Can study hours help predict exam score? & & \\[1.8em]
Can temperature help predict electricity demand? & & \\[1.8em]
Can a patient's symptoms help predict disease status? & & \\
\bottomrule
\end{tabularx}
\end{center}

\subsection*{Part 2: Model Assumptions}
The simple linear regression model is
\[
Y_i = \beta_0 + \beta_1x_i + \varepsilon_i.
\]
Match each assumption to a short interpretation.

\begin{center}
\renewcommand{\arraystretch}{1.45}
\begin{tabularx}{\textwidth}{p{4.3cm}X}
\toprule
\textbf{Assumption} & \textbf{What it means in context} \\
\midrule
\textbf{L}inearity & \\[2em]
\textbf{I}ndependent observations & \\[2em]
\textbf{N}ormal errors & \\[2em]
\textbf{E}qual (constant) variance & \\
\bottomrule
\end{tabularx}
\end{center}

\subsection*{Part 3: Rent and Apartment Size}
A rental website lists apartments in a large city. For each listing we record the size (square feet) and the monthly rent (dollars). A fitted linear model for predicting rent from size is
\[
\widehat{\text{rent}} = 800 + 2.5\times \text{size}.
\]

\begin{enumerate}[label={\bfseries (\Alph*)},itemsep=.75em]
\item What are the response variable and explanatory variable? \\[3em]
\item What relationship would you expect the scatterplot to show? Describe its likely direction, form, and strength. \\[3em]
\item What does the ``hat'' on $\widehat{\text{rent}}$ mean? \\[2.5em]
\item Predict the rent for a 700-square-foot apartment.
\[
\widehat{\text{rent}} = \blankline{3.5in}
\]
\item Interpret the slope $2.5$ in context. \\[3em]
\item The intercept is $800$. Is it practically meaningful here? Explain briefly. \\[3em]
\end{enumerate}

\subsection*{Part 4: Interpreting a Regression Model}
Suppose the fitted model is
\[
\widehat{\text{exam score}} = 52 + 4.8 \times \text{study hours}.
\]

\begin{enumerate}[label={\bfseries (\Alph*)},itemsep=.75em]
\item Interpret the intercept. Is it practically meaningful here? \\[3.5em]
\item Interpret the slope. \\[3em]
\item Predict the exam score for a student who studied 6 hours.
\[
\widehat{\text{exam score}} = \blankline{3.5in}
\]
\end{enumerate}

\subsection*{Part 5: Residuals and Least Squares}
Suppose we observe three data points and use the fitted line $\hat y = 1+x$.

\begin{center}
\renewcommand{\arraystretch}{1.55}
\begin{tabular}{ccccc}
\toprule
$x_i$ & $y_i$ & $\hat y_i$ & $e_i = y_i - \hat y_i$ & $e_i^2$ \\
\midrule
1 & 2 & & & \\[1.7em]
2 & 3 & & & \\[1.7em]
3 & 5 & & & \\
\bottomrule
\end{tabular}
\end{center}

\begin{enumerate}[label={\bfseries (\Alph*)},itemsep=.75em]
\item Compute the residual sum of squares.
\[
RSS = \sum_{i=1}^n e_i^2 = \blankline{2.5in}
\]
\item Which observations have positive residuals? What does a positive residual mean visually? \\[3em]
\end{enumerate}

\subsection*{Wrap-Up}
In one sentence, explain the connection between least squares and maximum likelihood under normal errors. \\[4em]

\end{document}
